% Balanced complete-edition layout; entries unchanged from the accepted bibliography.
\begin{thebibliography}{99}
\label{refs:chapter01}
\bibitem[Borgatti(2005)]{borgatti2005}Borgatti, S.~P. (2005).
Centrality and network flow.
\emph{Social Networks}, 27(1), 55--71.
\doi{10.1016/j.socnet.2004.11.008}.

\bibitem[Campbell(2011)]{campbell2011}Campbell, D.~T. (2011).
Assessing the impact of planned social change.
\emph{Journal of MultiDisciplinary Evaluation}, 7(15), 3--43.
\doi{10.56645/jmde.v7i15.297}.
Reprint of Occasional Paper Series No.~8, December 1976.

\bibitem[Crowston et~al.(2012)]{crowston2012}Crowston, K., Wei, K., Howison, J., and Wiggins, A. (2012).
Free/Libre open-source software development: What we know and what we do not know.
\emph{ACM Computing Surveys}, 44(2), Article~7.
\doi{10.1145/2089125.2089127}.

\bibitem[Dabbish et~al.(2012)]{dabbish2012}Dabbish, L., Stuart, C., Tsay, J., and Herbsleb, J. (2012).
Social coding in GitHub: Transparency and collaboration in an open software repository.
In \emph{Proceedings of the ACM 2012 Conference on Computer Supported Cooperative Work}, pp.~1277--1286. ACM.
\doi{10.1145/2145204.2145396}.

\bibitem[Decan et~al.(2019)]{decan2019}Decan, A., Mens, T., and Grosjean, P. (2019).
An empirical comparison of dependency network evolution in seven software packaging ecosystems.
\emph{Empirical Software Engineering}, 24(1), 381--416.
\doi{10.1007/s10664-017-9589-y}.

\bibitem[Fang and Neufeld(2009)]{fang2009}Fang, Y., and Neufeld, D. (2009).
Understanding sustained participation in open source software projects.
\emph{Journal of Management Information Systems}, 25(4), 9--50.
\doi{10.2753/MIS0742-1222250401}.

\bibitem[Forsgren et~al.(2021)]{forsgren2021}Forsgren, N., Storey, M.-A., Maddila, C., Zimmermann, T., Houck, B., and Butler, J. (2021).
The SPACE of developer productivity: There's more to it than you think.
\emph{ACM Queue}, 19(1), 20--48.
\doi{10.1145/3454122.3454124}.

\bibitem[Gousios et~al.(2014)]{gousios2014}Gousios, G., Pinzger, M., and van Deursen, A. (2014).
An exploratory study of the pull-based software development model.
In \emph{Proceedings of the 36th International Conference on Software Engineering}, pp.~345--355. ACM.
\doi{10.1145/2568225.2568260}.

\bibitem[Han et~al.(2024)]{han2024ose}Han, F., Zhao, S., Wang, W., Zhou, A., Qian, W., Zhou, X., Peng, J., You, L., Chen, Y., Xia, X., Tang, Y., Yang, L., and Tian, C. (2024).
Open Source Evaluatology: An evaluation framework and methodology for open source ecosystems based on evaluatology.
\emph{BenchCouncil Transactions on Benchmarks, Standards and Evaluations}, 4(4), 100190.
\doi{10.1016/j.tbench.2025.100190}.
\emph{Bibliographic note: the journal assigns this article to the December 2024 issue; the DOI and online publication date are in 2025.}

\bibitem[Han et~al.(2026a)]{han2026rank}Han, F., Zhao, S., Wang, W., Peng, J., and Xia, X. (2026a).
OpenRank: A centrality algorithm for high-dimensional heterogeneous networks in open source collaboration.
\emph{Information Sciences}, 730, 122893.
\doi{10.1016/j.ins.2025.122893}.

\bibitem[Howison and Crowston(2014)]{howison2014}Howison, J., and Crowston, K. (2014).
Collaboration through open superposition: A theory of the open source way.
\emph{MIS Quarterly}, 38(1), 29--50.
\doi{10.25300/MISQ/2014/38.1.02}.

\bibitem[Jansen(2014)]{jansen2014}Jansen, S. (2014).
Measuring the health of open source software ecosystems: Beyond the scope of project health.
\emph{Information and Software Technology}, 56(11), 1508--1519.
\doi{10.1016/j.infsof.2014.04.006}.

\bibitem[Kitchenham et~al.(1995)]{kitchenham1995}Kitchenham, B., Pfleeger, S.~L., and Fenton, N. (1995).
Towards a framework for software measurement validation.
\emph{IEEE Transactions on Software Engineering}, 21(12), 929--943.
\doi{10.1109/32.489070}.

\bibitem[Messick(1995)]{messick1995}Messick, S. (1995).
Validity of psychological assessment: Validation of inferences from persons' responses and performances as scientific inquiry into score meaning.
\emph{American Psychologist}, 50(9), 741--749.
\doi{10.1037/0003-066X.50.9.741}.

\end{thebibliography}
\clearpage
\begingroup
\renewcommand{\bibsection}{\section*{References\enspace{\normalfont\sffamily\fontsize{11}{14}\selectfont continued}}}
\begin{thebibliography}{99}
\bibitem[Mockus et~al.(2002)]{mockus2002}Mockus, A., Fielding, R.~T., and Herbsleb, J.~D. (2002).
Two case studies of open source software development: Apache and Mozilla.
\emph{ACM Transactions on Software Engineering and Methodology}, 11(3), 309--346.
\doi{10.1145/567793.567795}.

\bibitem[Shmueli(2010)]{shmueli2010}Shmueli, G. (2010).
To explain or to predict?
\emph{Statistical Science}, 25(3), 289--310.
\doi{10.1214/10-STS330}.

\bibitem[Steinmacher et~al.(2015)]{steinmacher2015}Steinmacher, I., Graciotto Silva, M.~A., Gerosa, M.~A., and Redmiles, D.~F. (2015).
A systematic literature review on the barriers faced by newcomers to open source software projects.
\emph{Information and Software Technology}, 59, 67--85.
\doi{10.1016/j.infsof.2014.11.001}.

\bibitem[Tsay et~al.(2014)]{tsay2014}Tsay, J., Dabbish, L., and Herbsleb, J. (2014).
Influence of social and technical factors for evaluating contribution in GitHub.
In \emph{Proceedings of the 36th International Conference on Software Engineering}, pp.~356--366. ACM.
\doi{10.1145/2568225.2568315}.

\bibitem[von Hippel and von Krogh(2003)]{vonhippel2003}von Hippel, E., and von Krogh, G. (2003).
Open source software and the “private-collective” innovation model: Issues for organization science.
\emph{Organization Science}, 14(2), 209--223.
\doi{10.1287/orsc.14.2.209.14992}.

\bibitem[Wang et~al.(2026)]{wang2026digger}Wang, W., Han, F., Zhao, S., Zhou, X., Qian, W., Zhou, A., Xia, X., Yang, L., Wang, R., Jiang, N., and Duan, M. (2026).
OpenDigger: A practical framework for assessing community health and sustainability in open source collaboration platforms.
In \emph{Proceedings of the ACM Web Conference 2026}, pp. 7745--7756. ACM.
\doi{10.1145/3774904.3792807}.

\bibitem[West and Gallagher(2006)]{west2006}West, J., and Gallagher, S. (2006).
Challenges of open innovation: The paradox of firm investment in open-source software.
\emph{R\&D Management}, 36(3), 319--331.
\doi{10.1111/j.1467-9310.2006.00436.x}.

\bibitem[Xia et~al.(2022)]{xia2022}Xia, X., Weng, Z., Wang, W., and Zhao, S. (2022).
Exploring activity and contributors on GitHub: Who, what, when, and where.
In \emph{2022 29th Asia-Pacific Software Engineering Conference}, pp. 11--20. IEEE.
\doi{10.1109/APSEC57359.2022.00013}.

\bibitem[Xia(2024)]{xia2024thesis}Xia, X. (2024).
\emph{Data-Driven Empirical Research on the Key Elements of Sustainability in Open Source Software Communities}.
Doctoral dissertation, East China Normal University, March 2024. In Chinese; English title supplied by the dissertation.

\bibitem[Zhan et~al.(2024)]{zhan2024}Zhan, J., Wang, L., Gao, W., Li, H., Wang, C., Huang, Y., Li, Y., Yang, Z., Kang, G., Luo, C., Ye, H., Dai, S., and Zhang, Z. (2024).
Evaluatology: The science and engineering of evaluation.
\emph{BenchCouncil Transactions on Benchmarks, Standards and Evaluations}, 4(1), 100162.
\doi{10.1016/j.tbench.2024.100162}.

\bibitem[Zhan et~al.(2026)]{zhan2026book}Zhan, J., Wang, L., Gao, W., Wang, C., Fan, F., Kang, G., and Li, H. (2026).
\emph{Evaluatology: The Science of Revealing Causes and Effects} [title translated from Chinese].
Chinese edition, version 1.0, August 2026. BenchCouncil Press.
ISBN 978-988-71596-9-8. \doi{10.66834/33wi1xx8}.
\emph{Page references refer to the printed pagination of this edition.}

\bibitem[Zhao et~al.(2024)]{zhao2024}Zhao, S., Xia, X., Fitzgerald, B., Li, X., Lenarduzzi, V., Taibi, D., Wang, R., Wang, W., and Tian, C. (2024).
OpenRank Leaderboard: Motivating open source collaborations through social network evaluation in Alibaba.
In \emph{Proceedings of the 46th International Conference on Software Engineering: Software Engineering in Practice}, pp. 346--357. ACM.
\doi{10.1145/3639477.3639734}.

\bibitem[Zhao(2025)]{zhao2025thesis}Zhao, S. (2025).
\emph{Research and Application of Key Technologies for Open Source Ecosystem Data Platform}.
Doctoral dissertation, Tongji University, September 2025. In Chinese; English title supplied by the dissertation.
\end{thebibliography}
\endgroup
